\documentclass[12pt,a4paper]{article}

% ── Page & fonts ─────────────────────────────────────────────────────────────
\usepackage[a4paper, top=2.5cm, bottom=2.5cm, left=2.8cm, right=2.8cm]{geometry}
\usepackage[T1]{fontenc}
\usepackage{lmodern}
\usepackage{microtype}

% ── Graphics & colour ────────────────────────────────────────────────────────
\usepackage{graphicx}
\usepackage{xcolor}
\usepackage{colortbl}

% ── Tables ───────────────────────────────────────────────────────────────────
\usepackage{booktabs}
\usepackage{multirow}
\usepackage{makecell}
\usepackage{array}
\usepackage{tabularx}

% ── Figures ──────────────────────────────────────────────────────────────────
\usepackage{subcaption}
\usepackage{float}
\usepackage{caption}
\captionsetup{font=small, labelfont=bf, skip=6pt}
\captionsetup[subfigure]{font=small, labelfont=bf}

% ── Boxes ────────────────────────────────────────────────────────────────────
\usepackage[most]{tcolorbox}
\usepackage{mdframed}

% ── Typography & layout ──────────────────────────────────────────────────────
\usepackage{titlesec}
\usepackage{parskip}
\usepackage{enumitem}
\usepackage{amsmath}
\usepackage{hyperref}
\usepackage{listings}
\usepackage{fancyhdr}

% ── Colour palette ───────────────────────────────────────────────────────────
\definecolor{primary}  {RGB}{31,  73, 125}   % dark blue
\definecolor{secondary}{RGB}{68, 114, 196}   % mid blue
\definecolor{accent}   {RGB}{237,125,  49}   % orange
\definecolor{lightbg}  {RGB}{242,242,242}    % light grey
\definecolor{wingreen} {RGB}{169,208,142}    % green
\definecolor{losered}  {RGB}{255,180,180}    % red
\definecolor{tieyellow}{RGB}{255,235,156}    % yellow
\definecolor{darktext} {RGB}{43, 43, 43}

% ── Table helpers ────────────────────────────────────────────────────────────
\newcommand{\win}[1]  {\cellcolor{wingreen}\textbf{#1}}
\newcommand{\lose}[1] {\cellcolor{losered}#1}
\newcommand{\tiew}[1] {\cellcolor{tieyellow}#1}
\newcommand{\hdr}[1]  {\cellcolor{secondary}\color{white}\textbf{#1}}
\newcommand{\hdrp}[1] {\cellcolor{primary}\color{white}\textbf{#1}}
\newcolumntype{C}[1]{>{\centering\arraybackslash}p{#1}}
\newcolumntype{L}[1]{>{\raggedright\arraybackslash}p{#1}}
\renewcommand{\arraystretch}{1.35}

% ── Section styling ──────────────────────────────────────────────────────────
\titleformat{\section}
  {\large\bfseries\color{primary}}
  {\thesection}{1em}{}
  [\color{secondary}\titlerule]

\titleformat{\subsection}
  {\normalsize\bfseries\color{secondary}}
  {\thesubsection}{1em}{}

% ── tcolorbox styles ─────────────────────────────────────────────────────────
\tcbset{
  summarybox/.style={
    colback=primary!8, colframe=primary, arc=4pt,
    boxrule=1pt, left=8pt, right=8pt, top=6pt, bottom=6pt,
    fonttitle=\bfseries\color{white}, coltitle=white,
    attach boxed title to top left={yshift=-2mm, xshift=6mm},
    boxed title style={colback=primary, arc=3pt}
  },
  findingbox/.style={
    colback=accent!10, colframe=accent, arc=4pt,
    boxrule=1pt, left=8pt, right=8pt, top=6pt, bottom=6pt,
    fonttitle=\bfseries\color{white}, coltitle=white,
    attach boxed title to top left={yshift=-2mm, xshift=6mm},
    boxed title style={colback=accent, arc=3pt}
  },
  codebox/.style={
    colback=lightbg, colframe=darktext!40, arc=3pt,
    boxrule=0.6pt, left=6pt, right=6pt, top=4pt, bottom=4pt,
    fontupper=\ttfamily\small
  },
  efftbox/.style={
    colback=wingreen!30, colframe=wingreen!80!black, arc=4pt,
    boxrule=1pt, left=6pt, right=6pt, top=4pt, bottom=4pt,
    fonttitle=\bfseries, title=EfficientNet-B0
  },
  resbox/.style={
    colback=secondary!15, colframe=secondary, arc=4pt,
    boxrule=1pt, left=6pt, right=6pt, top=4pt, bottom=4pt,
    fonttitle=\bfseries, title=ResNet50V2
  }
}

% ── Code listing style ───────────────────────────────────────────────────────
\lstset{
  language=Python,
  basicstyle=\ttfamily\footnotesize,
  keywordstyle=\color{secondary}\bfseries,
  stringstyle=\color{accent},
  commentstyle=\color{darktext!60}\itshape,
  backgroundcolor=\color{lightbg},
  frame=single, framerule=0.5pt, frameround=tttt,
  rulecolor=\color{darktext!30},
  breaklines=true, showstringspaces=false,
  numbers=none, xleftmargin=8pt, xrightmargin=8pt,
  aboveskip=4pt, belowskip=4pt
}

% ── Header / footer ──────────────────────────────────────────────────────────
\pagestyle{fancy}
\fancyhf{}
\fancyhead[L]{\small\color{darktext!60}Driver Drowsiness Detection — Model Comparison}
\fancyhead[R]{\small\color{darktext!60}Bachelor Thesis, April 2026}
\fancyfoot[C]{\small\thepage}
\renewcommand{\headrulewidth}{0.4pt}

% ─────────────────────────────────────────────────────────────────────────────
\begin{document}
% ─────────────────────────────────────────────────────────────────────────────

% ── Title page ───────────────────────────────────────────────────────────────
\begin{titlepage}
  \centering
  \vspace*{2cm}
  {\color{primary}\rule{\linewidth}{2pt}}\\[1em]
  {\Huge\bfseries\color{primary} Driver Drowsiness Detection}\\[0.6em]
  {\LARGE\color{secondary} EfficientNet-B0 vs.\ ResNet50V2}\\[0.4em]
  {\large Transfer Learning — Full Results Comparison}\\[0.5em]
  {\color{primary}\rule{\linewidth}{2pt}}\\[2em]

  \begin{tcolorbox}[summarybox, title=Quick Verdict, width=0.82\linewidth]
    \textbf{ResNet50V2 wins} on every meaningful accuracy metric:
    fewer errors (15 vs.\ 23), higher macro F1 (0.98 vs.\ 0.97),
    lower val loss (0.026 vs.\ 0.036), faster convergence (41 vs.\ 53 epochs),
    and cleaner feature embeddings.
    EfficientNet's only advantage is its $6\times$ smaller parameter count,
    relevant for edge/embedded deployment.
  \end{tcolorbox}

  \vspace{2em}
  {\large Bachelor Thesis\\[0.3em] April 2026}
  \vfill
\end{titlepage}

\tableofcontents
\newpage

% ═══════════════════════════════════════════════════════════════════════════════
\section{Dataset \& Experimental Setup}
% ═══════════════════════════════════════════════════════════════════════════════

Both models are trained and evaluated on the \textbf{same} 6-class drowsiness dataset
using an identical 70/30 random split (seed = 42).

\subsection{Classes}

\begin{center}
\begin{tabular}{clC{2.5cm}C{2.5cm}C{2.5cm}}
\toprule
\hdr{Label} & \hdr{Meaning} & \hdr{Train} & \hdr{Test} & \hdr{Total} \\
\midrule
0 & yawn        & $\sim$175  & 84  & $\sim$259 \\
1 & no\_yawn    & $\sim$160  & 76  & $\sim$236 \\
2 & Closed      & $\sim$500  & 223 & $\sim$723 \\
3 & Open        & $\sim$500  & 228 & $\sim$728 \\
4 & front       & $\sim$1430 & 594 & $\sim$2024 \\
5 & down        & $\sim$1540 & 642 & $\sim$2182 \\
\midrule
  & \textbf{Total} & \textbf{4{,}313} & \textbf{1{,}847} & \textbf{6{,}160} \\
\bottomrule
\end{tabular}
\end{center}

\begin{tcolorbox}[findingbox, title=Dataset Imbalance Note]
\textit{front} and \textit{down} have roughly $8\times$ more samples than
\textit{yawn} and \textit{no\_yawn}.
This directly explains why both models struggle most on \textit{yawn} — it has the
smallest support (84 test samples) and is the hardest class visually.
\end{tcolorbox}

\begin{figure}[H]
  \centering
  \begin{subfigure}[t]{0.48\textwidth}
    \includegraphics[width=\linewidth]{images/efficientnet_class_dist.png}
    \caption{EfficientNet-B0}
  \end{subfigure}
  \hfill
  \begin{subfigure}[t]{0.48\textwidth}
    \includegraphics[width=\linewidth]{images/resnet_class_dist.png}
    \caption{ResNet50V2}
  \end{subfigure}
  \caption{Class distribution — Train vs.\ Test split (identical for both models)}
\end{figure}

% ═══════════════════════════════════════════════════════════════════════════════
\section{Code Architecture \& Logic Differences}
% ═══════════════════════════════════════════════════════════════════════════════

The two notebooks share the same data-loading pipeline and two-phase training
strategy, but differ in base model, preprocessing, and head architecture.

% ──────────────────────────────────────────────────────────────────────────────
\subsection{Data Loading Pipeline (Shared)}

Both notebooks load data through three separate functions and merge them:

\begin{lstlisting}
# Step 1 — yawn / no_yawn: Haar-cascade face crop before resizing
def face_for_yawn(direc, face_cas_path):
    face_cascade = cv2.CascadeClassifier(face_cas_path)
    faces = face_cascade.detectMultiScale(image_array, 1.3, 5)
    roi_color = img[y:y+h, x:x+w]          # crop to face ROI
    resized = cv2.resize(roi_color, (224, 224))

# Step 2 — Closed / Open eyes: direct resize, no face crop
def get_data(dir_path):
    resized = cv2.resize(img_array, (224, 224))

# Step 3 — head pose front / down: direct resize with validity checks
def get_head_pose_data(dir_path):
    if img_array is None: continue
    if resized.shape != (224, 224, 3): continue
\end{lstlisting}

\begin{tcolorbox}[findingbox, title=Important Design Decision]
\textit{yawn} and \textit{no\_yawn} go through a Haar-cascade face detector before
resizing, while \textit{Closed}/\textit{Open} eyes and head-pose images are resized
directly. This means yawn/no\_yawn images are face-cropped patches, whereas other
classes are full-frame crops — a potential source of domain mismatch.
\end{tcolorbox}

% ──────────────────────────────────────────────────────────────────────────────
\subsection{Preprocessing}

This is the most critical code difference between the two notebooks:

\begin{tcolorbox}[efftbox]
\begin{lstlisting}
# EfficientNet-B0 — NO explicit preprocessing needed.
# The model contains an internal rescaling layer.
train_gen = ImageDataGenerator(
    zoom_range=0.2, horizontal_flip=True, rotation_range=30
    # no rescale, no preprocessing_function
)
# At inference time: pass raw uint8 images (0-255)
y_prob = model.predict(X_test.astype('float32'))
\end{lstlisting}
\end{tcolorbox}

\begin{tcolorbox}[resbox]
\begin{lstlisting}
from tensorflow.keras.applications.resnet_v2 import preprocess_input
# ResNet50V2 — MUST apply preprocess_input explicitly.
# Scales pixels from [0,255] to [-1,1].
train_gen = ImageDataGenerator(
    preprocessing_function=preprocess_input,   # <-- required
    zoom_range=0.2, horizontal_flip=True, rotation_range=30
)
# At inference time: always preprocess first
X_test_proc = preprocess_input(X_test.astype('float32'))
y_prob = model.predict(X_test_proc)
\end{lstlisting}
\end{tcolorbox}

\begin{tcolorbox}[findingbox, title=Why This Matters]
Forgetting \texttt{preprocess\_input} with ResNet50V2 would feed pixel values in
[0, 255] to a network that expects [-1, 1], causing dramatically wrong predictions.
EfficientNet avoids this pitfall with its built-in rescaling layer.
The same preprocessing must also be applied to Grad-CAM inputs and the
t-SNE feature extractor — three places where the two notebooks diverge.
\end{tcolorbox}

% ──────────────────────────────────────────────────────────────────────────────
\subsection{Model Head Architecture}

\begin{tcolorbox}[efftbox]
\begin{lstlisting}
base_model = EfficientNetB0(weights='imagenet', include_top=False,
                            input_shape=(224,224,3))
base_model.trainable = False

x = GlobalAveragePooling2D()(base_model.output)   # 1280-dim
x = Dense(256, activation='relu')(x)
x = Dropout(0.5)(x)
output = Dense(6, activation='softmax')(x)         # 3 layers total
\end{lstlisting}
\end{tcolorbox}

\begin{tcolorbox}[resbox]
\begin{lstlisting}
base_model = ResNet50V2(weights='imagenet', include_top=False,
                        input_shape=(224,224,3))
base_model.trainable = False

x = GlobalAveragePooling2D()(base_model.output)   # 2048-dim
x = BatchNormalization()(x)     # extra BN — matches ResNet's design
x = Dense(256, activation='relu')(x)
x = Dropout(0.5)(x)
x = Dense(128, activation='relu')(x)              # extra hidden layer
x = Dropout(0.3)(x)
output = Dense(6, activation='softmax')(x)         # 5 layers total
\end{lstlisting}
\end{tcolorbox}

\begin{table}[H]
\centering
\caption{Head architecture comparison}
\begin{tabular}{L{4.5cm}C{4cm}C{4cm}}
\toprule
\hdr{Property} & \hdr{EfficientNet-B0} & \hdr{ResNet50V2} \\
\midrule
GAP output dim      & 1{,}280          & 2{,}048 \\
BatchNorm after GAP & No               & \textbf{Yes} \\
Hidden layers       & 1 $\times$ Dense(256) & 2 $\times$ Dense(256, 128) \\
Dropout             & 0.5              & 0.5 + 0.3 \\
Total head params   & $\sim$330K       & $\sim$590K \\
\bottomrule
\end{tabular}
\end{table}

% ──────────────────────────────────────────────────────────────────────────────
\subsection{Two-Phase Training Strategy (Shared Logic)}

\begin{lstlisting}
# ── Phase 1: freeze base, train head only ──────────────────────────────────
base_model.trainable = False
model.compile(optimizer=Adam(lr=1e-3), loss='categorical_crossentropy')
history1 = model.fit(..., epochs=20,
    callbacks=[EarlyStopping(patience=5), ReduceLROnPlateau(patience=3)])

# ── Phase 2: unfreeze last 30 layers, fine-tune ────────────────────────────
base_model.trainable = True
for layer in base_model.layers[:-30]:
    layer.trainable = False
model.compile(optimizer=Adam(lr=1e-5), ...)   # 100x lower LR
history2 = model.fit(..., epochs=40,
    callbacks=[EarlyStopping(patience=7), ReduceLROnPlateau(patience=3)])
\end{lstlisting}

% ═══════════════════════════════════════════════════════════════════════════════
\section{Training Results}
% ═══════════════════════════════════════════════════════════════════════════════

\subsection{Training Curves}

\begin{figure}[H]
  \centering
  \begin{subfigure}[t]{0.48\textwidth}
    \includegraphics[width=\linewidth]{images/efficientnet_accuracy.png}
    \caption{EfficientNet-B0 — Accuracy}
  \end{subfigure}
  \hfill
  \begin{subfigure}[t]{0.48\textwidth}
    \includegraphics[width=\linewidth]{images/resnet_accuracy.png}
    \caption{ResNet50V2 — Accuracy}
  \end{subfigure}

  \vspace{0.6em}

  \begin{subfigure}[t]{0.48\textwidth}
    \includegraphics[width=\linewidth]{images/efficientnet_loss.png}
    \caption{EfficientNet-B0 — Loss}
  \end{subfigure}
  \hfill
  \begin{subfigure}[t]{0.48\textwidth}
    \includegraphics[width=\linewidth]{images/resnet_loss.png}
    \caption{ResNet50V2 — Loss}
  \end{subfigure}
  \caption{Training and validation accuracy/loss over all epochs.
           Dashed vertical line marks the Phase 1 $\to$ Phase 2 boundary.
           The large spike in training loss at the boundary is expected
           when unfreezing layers.}
\end{figure}

\subsection{Phase-by-Phase Training Summary}

\begin{table}[H]
\centering
\caption{Phase 1 — Frozen base, head training only}
\begin{tabular}{L{5cm}C{3.8cm}C{3.8cm}}
\toprule
\hdr{Metric} & \hdr{EfficientNet-B0} & \hdr{ResNet50V2} \\
\midrule
Starting LR                   & $10^{-3}$             & $10^{-3}$ \\
Ep 1 train / val acc          & 84.0\% / 96.7\%       & 79.7\% / 95.2\% \\
Best val accuracy             & 98.38\% (ep 13)       & \win{98.48\% (ep 15)} \\
Best val loss                 & \lose{0.0428 (ep 8)}  & \win{0.0482 (ep 14)} \\
Epochs run / max              & \lose{13 / 20}        & 19 / 20 \\
Early stopped                 & Yes (patience 5)      & Yes (patience 5) \\
Final LR                      & $1.25 \times 10^{-4}$ & $1.25 \times 10^{-4}$ \\
\bottomrule
\end{tabular}
\end{table}

\begin{table}[H]
\centering
\caption{Phase 2 — Partial fine-tuning (last 30 layers unfrozen)}
\begin{tabular}{L{5cm}C{3.8cm}C{3.8cm}}
\toprule
\hdr{Metric} & \hdr{EfficientNet-B0} & \hdr{ResNet50V2} \\
\midrule
Starting LR                    & $10^{-5}$              & $10^{-5}$ \\
Ep 1 train / val acc           & 86.1\% / 98.4\%        & 96.7\% / 98.4\% \\
Best val accuracy              & 98.81\%                & \win{99.19\%} \\
Best val loss                  & \lose{0.0360 (ep 40)}  & \win{0.0258 (ep 15)} \\
Epochs run / max               & \lose{40 / 40 (all)}   & \win{22 / 40} \\
Early stopped                  & \lose{No — never converged} & \win{Yes (patience 7)} \\
Final LR                       & $6.25 \times 10^{-7}$  & $1.25 \times 10^{-6}$ \\
\bottomrule
\end{tabular}
\end{table}

\begin{table}[H]
\centering
\caption{Overall training comparison}
\begin{tabular}{L{5cm}C{3.8cm}C{3.8cm}C{2cm}}
\toprule
\hdr{Metric} & \hdr{EfficientNet-B0} & \hdr{ResNet50V2} & \hdr{Winner} \\
\midrule
Total epochs           & \lose{53}         & \win{41}          & ResNet \\
Best val accuracy      & \lose{98.81\%}    & \win{99.19\%}     & ResNet \\
Best val loss          & \lose{0.0360}     & \win{0.0258}      & ResNet \\
Convergence behaviour  & Slow, noisy       & \win{Fast, smooth}& ResNet \\
\bottomrule
\end{tabular}
\end{table}

\begin{tcolorbox}[findingbox, title=Key Training Observation]
EfficientNet's Phase 2 ran all 40 epochs without triggering early stopping,
meaning it never truly plateaued — val loss kept slowly decreasing (0.0482 $\to$ 0.0360)
with no sign of overfitting. ResNet converged cleanly after just 22 epochs, achieving a
significantly lower val loss of 0.0258. This suggests EfficientNet could benefit from
even more Phase 2 epochs, while ResNet found a better minimum.
\end{tcolorbox}

% ═══════════════════════════════════════════════════════════════════════════════
\section{Overall Test Performance}
% ═══════════════════════════════════════════════════════════════════════════════

\begin{table}[H]
\centering
\caption{Overall metrics on the 1{,}847-sample test set}
\begin{tabular}{L{5cm}C{3.8cm}C{3.8cm}C{2cm}}
\toprule
\hdr{Metric} & \hdr{EfficientNet-B0} & \hdr{ResNet50V2} & \hdr{Winner} \\
\midrule
Overall accuracy       & \tiew{99\%}        & \tiew{99\%}       & Tie \\
Macro avg precision    & \lose{0.97}        & \win{0.98}        & ResNet \\
Macro avg recall       & \lose{0.97}        & \win{0.98}        & ResNet \\
Macro avg F1           & \lose{0.97}        & \win{0.98}        & ResNet \\
Weighted avg precision & \tiew{0.99}        & \tiew{0.99}       & Tie \\
Weighted avg recall    & \tiew{0.99}        & \tiew{0.99}       & Tie \\
Weighted avg F1        & \tiew{0.99}        & \tiew{0.99}       & Tie \\
Total misclassified    & \lose{23 / 1{,}847}& \win{15 / 1{,}847}& ResNet \\
\bottomrule
\end{tabular}
\end{table}

% ═══════════════════════════════════════════════════════════════════════════════
\section{Per-Class Classification Report}
% ═══════════════════════════════════════════════════════════════════════════════

\begin{table}[H]
\centering
\caption{Full per-class classification report — Precision / Recall / F1 / Support}
\begin{tabular}{lC{1.7cm}C{1.7cm}C{1.7cm}C{1cm}C{1.7cm}C{1.7cm}C{1.7cm}C{1cm}}
\toprule
 & \multicolumn{4}{c}{\cellcolor{secondary}\color{white}\textbf{EfficientNet-B0}}
 & \multicolumn{4}{c}{\cellcolor{primary}\color{white}\textbf{ResNet50V2}} \\
\cmidrule(lr){2-5}\cmidrule(lr){6-9}
\textbf{Class} & Prec & Rec & F1 & Sup & Prec & Rec & F1 & Sup \\
\midrule
yawn        & 0.97        & \lose{0.86} & \lose{0.91} & 84
            & 0.97        & \win{0.89}  & \win{0.93}  & 84 \\
no\_yawn    & \lose{0.87} & 0.97        & \lose{0.92} & 76
            & \win{0.90}  & 0.97        & \win{0.94}  & 76 \\
Closed      & \lose{0.99} & 1.00        & \lose{0.99} & 223
            & \win{1.00}  & \win{1.00}  & \win{1.00}  & 223 \\
Open        & 1.00        & \lose{0.99} & 1.00        & 228
            & 1.00        & \win{1.00}  & 1.00        & 228 \\
front       & 1.00        & \lose{0.99} & \lose{0.99} & 594
            & 1.00        & \win{0.99}  & \win{1.00}  & 594 \\
down        & \lose{0.99} & 1.00        & \lose{0.99} & 642
            & \win{1.00}  & 1.00        & \win{1.00}  & 642 \\
\midrule
macro avg   & \lose{0.97} & \lose{0.97} & \lose{0.97} & 1847
            & \win{0.98}  & \win{0.98}  & \win{0.98}  & 1847 \\
wtd avg     & 0.99        & 0.99        & 0.99        & 1847
            & 0.99        & 0.99        & 0.99        & 1847 \\
\bottomrule
\end{tabular}
\end{table}

% ═══════════════════════════════════════════════════════════════════════════════
\section{Confusion Matrices}
% ═══════════════════════════════════════════════════════════════════════════════

\begin{figure}[H]
  \centering
  \begin{subfigure}[t]{0.48\textwidth}
    \includegraphics[width=\linewidth]{images/efficientnet_confusion_matrix.png}
    \caption{EfficientNet-B0 — \textbf{23 total errors}}
  \end{subfigure}
  \hfill
  \begin{subfigure}[t]{0.48\textwidth}
    \includegraphics[width=\linewidth]{images/resnet_confusion_matrix.png}
    \caption{ResNet50V2 — \textbf{15 total errors}}
  \end{subfigure}
  \caption{Confusion matrices on the 1{,}847-sample test set.
           Both models make errors almost exclusively in the
           \textit{yawn}/\textit{no\_yawn} and \textit{front}/\textit{down} pairs.}
\end{figure}

\begin{table}[H]
\centering
\caption{Exact confusion matrix values}
\begin{tabular}{lC{1.5cm}C{1.5cm}C{1.5cm}C{1.5cm}C{1.5cm}C{1.5cm}}
\toprule
\multicolumn{7}{c}{\cellcolor{secondary}\color{white}\textbf{EfficientNet-B0}} \\
\midrule
\textbf{True\textbackslash Pred} & yawn & no\_yn & Closed & Open & front & down \\
\midrule
yawn    & \win{72} & \lose{11} & 1 & 0 & 0 & 0 \\
no\_yawn& 2 & \win{74} & 0 & 0 & 0 & 0 \\
Closed  & 0 & 0 & \win{223} & 0 & 0 & 0 \\
Open    & 0 & 0 & \lose{2}  & \win{226} & 0 & 0 \\
front   & 0 & 0 & 0 & 0 & \win{587} & \lose{7} \\
down    & 0 & 0 & 0 & 0 & 0 & \win{642} \\
\bottomrule
\end{tabular}

\vspace{0.8em}

\begin{tabular}{lC{1.5cm}C{1.5cm}C{1.5cm}C{1.5cm}C{1.5cm}C{1.5cm}}
\toprule
\multicolumn{7}{c}{\cellcolor{primary}\color{white}\textbf{ResNet50V2}} \\
\midrule
\textbf{True\textbackslash Pred} & yawn & no\_yn & Closed & Open & front & down \\
\midrule
yawn    & \win{75} & \lose{8} & 1 & 0 & 0 & 0 \\
no\_yawn& 2 & \win{74} & 0 & 0 & 0 & 0 \\
Closed  & 0 & 0 & \win{222} & 1 & 0 & 0 \\
Open    & 0 & 0 & 0 & \win{228} & 0 & 0 \\
front   & 0 & 0 & 0 & 0 & \win{591} & \lose{3} \\
down    & 0 & 0 & 0 & 0 & 0 & \win{642} \\
\bottomrule
\end{tabular}
\end{table}

\begin{table}[H]
\centering
\caption{Error-by-error breakdown}
\begin{tabular}{L{4.5cm}C{3.5cm}C{3.5cm}C{2.2cm}}
\toprule
\hdr{Error type} & \hdr{EfficientNet-B0} & \hdr{ResNet50V2} & \hdr{Winner} \\
\midrule
yawn $\to$ no\_yawn & \lose{11}  & \win{8}  & ResNet \\
yawn $\to$ Closed   & 1          & 1        & Tie \\
no\_yawn $\to$ yawn & 2          & 2        & Tie \\
Closed $\to$ Open   & \win{0}    & \lose{1} & EfficientNet \\
Open $\to$ Closed   & \lose{2}   & \win{0}  & ResNet \\
front $\to$ down    & \lose{7}   & \win{3}  & ResNet \\
\midrule
\textbf{Total}      & \lose{\textbf{23}} & \win{\textbf{15}} & \textbf{ResNet} \\
\bottomrule
\end{tabular}
\end{table}

% ═══════════════════════════════════════════════════════════════════════════════
\section{Per-Class Accuracy}
% ═══════════════════════════════════════════════════════════════════════════════

\begin{figure}[H]
  \centering
  \begin{subfigure}[t]{0.48\textwidth}
    \includegraphics[width=\linewidth]{images/efficientnet_per_class_accuracy.png}
    \caption{EfficientNet-B0}
  \end{subfigure}
  \hfill
  \begin{subfigure}[t]{0.48\textwidth}
    \includegraphics[width=\linewidth]{images/resnet_per_class_accuracy.png}
    \caption{ResNet50V2}
  \end{subfigure}
  \caption{Per-class accuracy. The \textit{yawn} class is the weak point for both models,
           with ResNet achieving 89.3\% vs.\ EfficientNet's 85.7\%.}
\end{figure}

\begin{table}[H]
\centering
\caption{Per-class accuracy — exact values}
\begin{tabular}{lC{3.8cm}C{3.8cm}C{2.5cm}}
\toprule
\hdr{Class} & \hdr{EfficientNet-B0} & \hdr{ResNet50V2} & \hdr{Winner} \\
\midrule
yawn     & \lose{85.7\%} & \win{89.3\%}  & ResNet \\
no\_yawn & \tiew{97.4\%} & \tiew{97.4\%} & Tie \\
Closed   & \win{100.0\%} & \lose{99.6\%} & EfficientNet \\
Open     & \lose{99.1\%} & \win{100.0\%} & ResNet \\
front    & \lose{98.8\%} & \win{99.5\%}  & ResNet \\
down     & \tiew{100.0\%}& \tiew{100.0\%}& Tie \\
\bottomrule
\end{tabular}
\end{table}

% ═══════════════════════════════════════════════════════════════════════════════
\section{ROC Curves \& AUC Scores}
% ═══════════════════════════════════════════════════════════════════════════════

\begin{figure}[H]
  \centering
  \begin{subfigure}[t]{0.48\textwidth}
    \includegraphics[width=\linewidth]{images/efficientnet_roc.png}
    \caption{EfficientNet-B0}
  \end{subfigure}
  \hfill
  \begin{subfigure}[t]{0.48\textwidth}
    \includegraphics[width=\linewidth]{images/resnet_roc.png}
    \caption{ResNet50V2}
  \end{subfigure}
  \caption{ROC curves (one-vs-rest strategy). Both models achieve near-perfect AUC.
           All curves hug the top-left corner, indicating excellent class discrimination
           even for the weaker \textit{yawn}/\textit{no\_yawn} classes.}
\end{figure}

\begin{table}[H]
\centering
\caption{ROC AUC per class — one-vs-rest}
\begin{tabular}{lC{3.8cm}C{3.8cm}C{2.5cm}}
\toprule
\hdr{Class} & \hdr{EfficientNet-B0} & \hdr{ResNet50V2} & \hdr{Winner} \\
\midrule
yawn     & \tiew{0.999} & \tiew{0.999} & Tie \\
no\_yawn & \tiew{0.999} & \tiew{0.999} & Tie \\
Closed   & \tiew{1.000} & \tiew{1.000} & Tie \\
Open     & \tiew{1.000} & \tiew{1.000} & Tie \\
front    & \tiew{1.000} & \tiew{1.000} & Tie \\
down     & \tiew{1.000} & \tiew{1.000} & Tie \\
\bottomrule
\end{tabular}
\end{table}

\begin{tcolorbox}[summarybox, title=AUC Interpretation]
AUC is insensitive to the small differences between the models here — both achieve
essentially perfect discrimination at the ranking level.
The precision-recall curves and confusion matrices are more informative for this
near-perfect but imbalanced setting.
\end{tcolorbox}

% ═══════════════════════════════════════════════════════════════════════════════
\section{Precision-Recall Curves}
% ═══════════════════════════════════════════════════════════════════════════════

\begin{figure}[H]
  \centering
  \begin{subfigure}[t]{0.48\textwidth}
    \includegraphics[width=\linewidth]{images/efficientnet_pr.png}
    \caption{EfficientNet-B0}
  \end{subfigure}
  \hfill
  \begin{subfigure}[t]{0.48\textwidth}
    \includegraphics[width=\linewidth]{images/resnet_pr.png}
    \caption{ResNet50V2}
  \end{subfigure}
  \caption{Precision-Recall curves with Average Precision (AP) scores.
           For the harder \textit{yawn} and \textit{no\_yawn} classes,
           ResNet achieves slightly higher AP, reflecting better confidence calibration.}
\end{figure}

\begin{table}[H]
\centering
\caption{Average Precision (AP) per class}
\begin{tabular}{lC{3.8cm}C{3.8cm}C{2.5cm}}
\toprule
\hdr{Class} & \hdr{EfficientNet-B0} & \hdr{ResNet50V2} & \hdr{Winner} \\
\midrule
yawn     & \lose{0.981} & \win{0.986} & ResNet \\
no\_yawn & \lose{0.975} & \win{0.985} & ResNet \\
Closed   & \tiew{1.000} & \tiew{1.000} & Tie \\
Open     & \tiew{1.000} & \tiew{1.000} & Tie \\
front    & \tiew{1.000} & \tiew{1.000} & Tie \\
down     & \tiew{1.000} & \tiew{1.000} & Tie \\
\bottomrule
\end{tabular}
\end{table}

% ═══════════════════════════════════════════════════════════════════════════════
\section{Error Analysis — Misclassification Gallery}
% ═══════════════════════════════════════════════════════════════════════════════

\begin{figure}[H]
  \centering
  \begin{subfigure}[t]{0.48\textwidth}
    \includegraphics[width=\linewidth]{images/efficientnet_misclassified.png}
    \caption{EfficientNet-B0 — \textbf{23 misclassified} images}
  \end{subfigure}
  \hfill
  \begin{subfigure}[t]{0.48\textwidth}
    \includegraphics[width=\linewidth]{images/resnet_misclassified.png}
    \caption{ResNet50V2 — \textbf{15 misclassified} images}
  \end{subfigure}
  \caption{Up to 16 misclassified images shown per model (true label top, predicted label bottom).}
\end{figure}

\begin{tcolorbox}[findingbox, title=Error Pattern Analysis]
\begin{itemize}[noitemsep, leftmargin=1.2em]
  \item \textbf{yawn $\leftrightarrow$ no\_yawn}: The dominant error for both models.
        These classes are inherently ambiguous — subtle mouth openings without
        full yawn stretching are hard to distinguish. The Haar-cascade face crop
        used during loading may also discard context that helps disambiguate.
  \item \textbf{front $\to$ down}: EfficientNet makes 7 such errors vs.\ ResNet's 3.
        Both classes are head-pose images, and subtle downward tilts at the
        boundary of the two classes are visually similar.
  \item \textbf{Closed/Open}: Both models are near-perfect here;
        EfficientNet makes no Open$\to$Closed errors but 2 Open$\to$Closed errors,
        while ResNet makes only 1 Closed$\to$Open error.
\end{itemize}
\end{tcolorbox}

% ═══════════════════════════════════════════════════════════════════════════════
\section{Interpretability — Grad-CAM Heatmaps}
% ═══════════════════════════════════════════════════════════════════════════════

\begin{figure}[H]
  \centering
  \begin{subfigure}[t]{\textwidth}
    \includegraphics[width=\linewidth]{images/efficientnet_gradcam.png}
    \caption{EfficientNet-B0 — target layer: \texttt{top\_conv}\\
             \small Top row: original images. Bottom row: Grad-CAM overlay.}
  \end{subfigure}
  \vspace{0.8em}
  \begin{subfigure}[t]{\textwidth}
    \includegraphics[width=\linewidth]{images/resnet_gradcam.png}
    \caption{ResNet50V2 — target layer: \texttt{conv5\_block3\_3\_conv}\\
             \small Top row: original images. Bottom row: Grad-CAM overlay.}
  \end{subfigure}
  \caption{Grad-CAM heatmaps — one representative sample per class.
           The heatmap colour goes from blue (low attention) to red (high attention).}
\end{figure}

\begin{table}[H]
\centering
\caption{Grad-CAM attention quality — qualitative assessment}
\begin{tabular}{lL{4.5cm}L{4.5cm}}
\toprule
\hdr{Class} & \hdr{EfficientNet-B0} & \hdr{ResNet50V2} \\
\midrule
yawn / no\_yawn & Broad face + mouth region   & Tighter mouth focus \\
Closed / Open   & Eye region (diffuse spread) & Eye region (more compact) \\
front / down    & Head + some background      & Head orientation, less background \\
\midrule
\textbf{Overall} & Correct but diffuse        & \textbf{More precise and localised} \\
\bottomrule
\end{tabular}
\end{table}

\begin{tcolorbox}[summarybox, title=Grad-CAM Significance for the Thesis]
Both models attend to the correct facial regions — mouth for yawn, eyes for
closed/open, head orientation for front/down. This confirms the models are not
exploiting background shortcuts or dataset artefacts.
ResNet's tighter activations suggest its deeper head (with extra BatchNorm and
hidden layer) produces more discriminative, spatially-focused features.
\end{tcolorbox}

% ═══════════════════════════════════════════════════════════════════════════════
\section{Interpretability — t-SNE Feature Embeddings}
% ═══════════════════════════════════════════════════════════════════════════════

\begin{figure}[H]
  \centering
  \begin{subfigure}[t]{0.48\textwidth}
    \includegraphics[width=\linewidth]{images/efficientnet_tsne.png}
    \caption{EfficientNet-B0 — more class overlap}
  \end{subfigure}
  \hfill
  \begin{subfigure}[t]{0.48\textwidth}
    \includegraphics[width=\linewidth]{images/resnet_tsne.png}
    \caption{ResNet50V2 — better class separation}
  \end{subfigure}
  \caption{t-SNE projection of GlobalAveragePooling2D layer embeddings
           (up to 2{,}000 randomly sampled test images).
           Tighter, more separated clusters indicate richer learned representations.}
\end{figure}

\begin{table}[H]
\centering
\caption{t-SNE cluster quality — qualitative comparison}
\begin{tabular}{lL{4.5cm}L{4.5cm}}
\toprule
\hdr{Class} & \hdr{EfficientNet-B0} & \hdr{ResNet50V2} \\
\midrule
Closed   & Partially mixed with Open/front & \textbf{Tight, clean cluster} \\
Open     & Partially mixed with Closed     & \textbf{Tight, clean cluster} \\
yawn     & Corner cluster, some overlap    & Overlaps with no\_yawn \\
no\_yawn & Mixed with yawn                 & Overlaps with yawn \\
front    & Very fragmented across plot     & Fragmented, more contained \\
down     & Scattered sub-clusters          & Clear distinct sub-clusters \\
\midrule
\textbf{Overall} & More class overlap; less separable & \textbf{Better separation overall} \\
\bottomrule
\end{tabular}
\end{table}

\begin{tcolorbox}[findingbox, title=t-SNE Interpretation]
The yawn/no\_yawn overlap in \textbf{both} models' embedding spaces directly
explains why those two classes have the worst classification scores —
the network representations themselves are not fully separable for these classes,
which is consistent with the confusion matrix errors.
ResNet's better-separated Closed/Open clusters explain its perfect 100\% accuracy
on those two classes vs.\ EfficientNet's 99.6\% / 99.1\%.
\end{tcolorbox}

% ═══════════════════════════════════════════════════════════════════════════════
\section{Final Verdict}
% ═══════════════════════════════════════════════════════════════════════════════

\begin{table}[H]
\centering
\caption{Head-to-head summary across all evaluation dimensions}
\begin{tabular}{L{5.5cm}C{3.2cm}C{3.2cm}C{2cm}}
\toprule
\hdr{Dimension} & \hdr{EfficientNet-B0} & \hdr{ResNet50V2} & \hdr{Winner} \\
\midrule
\multicolumn{4}{l}{\cellcolor{lightbg}\textit{Training}} \\
Total epochs          & \lose{53}          & \win{41}          & ResNet \\
Best val accuracy     & \lose{98.81\%}     & \win{99.19\%}     & ResNet \\
Best val loss         & \lose{0.0360}      & \win{0.0258}      & ResNet \\
Convergence           & \lose{Slow, noisy} & \win{Fast, smooth}& ResNet \\
\midrule
\multicolumn{4}{l}{\cellcolor{lightbg}\textit{Test Performance}} \\
Overall accuracy      & \tiew{99\%}        & \tiew{99\%}       & Tie \\
Macro F1              & \lose{0.97}        & \win{0.98}        & ResNet \\
Total errors          & \lose{23}          & \win{15}          & ResNet \\
Yawn accuracy         & \lose{85.7\%}      & \win{89.3\%}      & ResNet \\
No\_yawn accuracy     & \tiew{97.4\%}      & \tiew{97.4\%}     & Tie \\
Closed accuracy       & \win{100.0\%}      & \lose{99.6\%}     & EfficientNet \\
Open accuracy         & \lose{99.1\%}      & \win{100.0\%}     & ResNet \\
Front accuracy        & \lose{98.8\%}      & \win{99.5\%}      & ResNet \\
Down accuracy         & \tiew{100.0\%}     & \tiew{100.0\%}    & Tie \\
\midrule
\multicolumn{4}{l}{\cellcolor{lightbg}\textit{Ranking Metrics}} \\
AUC (yawn/no\_yawn)   & \tiew{0.999}       & \tiew{0.999}      & Tie \\
AP (yawn)             & \lose{0.981}       & \win{0.986}       & ResNet \\
AP (no\_yawn)         & \lose{0.975}       & \win{0.985}       & ResNet \\
\midrule
\multicolumn{4}{l}{\cellcolor{lightbg}\textit{Interpretability}} \\
Grad-CAM localisation & Correct, diffuse   & \win{Correct, precise} & ResNet \\
t-SNE separability    & \lose{More overlap}& \win{Better separation}& ResNet \\
\midrule
\multicolumn{4}{l}{\cellcolor{lightbg}\textit{Practical Considerations}} \\
Parameters            & \win{$\sim$4M}    & \lose{$\sim$23.5M}& EfficientNet \\
Inference speed       & \win{Faster}      & \lose{Slower}     & EfficientNet \\
Edge deployment       & \win{Suitable}    & \lose{Heavier}    & EfficientNet \\
\midrule
\textbf{Total wins}   & \lose{\textbf{2}} & \win{\textbf{12}} & \textbf{ResNet} \\
\bottomrule
\end{tabular}
\end{table}

\vspace{1em}

\begin{tcolorbox}[summarybox, title=Conclusion]
\textbf{ResNet50V2 is the stronger model} for this drowsiness detection task,
winning on 12 out of 14 dimensions evaluated.
It produces 35\% fewer errors (15 vs.\ 23), better macro F1 (0.98 vs.\ 0.97),
lower validation loss (0.026 vs.\ 0.036), faster and smoother convergence
(41 vs.\ 53 epochs), and more separable internal feature representations.

\medskip
The \textbf{yawn class remains the hardest} for both models (85.7\% vs.\ 89.3\%),
caused by a combination of small sample size (84 test samples), inherent visual
ambiguity with no\_yawn, and the Haar-cascade face-crop pipeline discarding
context that might help disambiguate.

\medskip
EfficientNet's sole advantages are its \textbf{$6\times$ smaller parameter count}
($\sim$4M vs.\ $\sim$23.5M) and faster inference — relevant if the model
is deployed on embedded hardware (e.g., a Raspberry Pi in a vehicle).
For server-side or GPU inference where accuracy is the priority, ResNet50V2
is the clear choice.
\end{tcolorbox}

\end{document}
